\section{Pulp Fiction}

\begin{quotex}
There are few people who know the full force of the different movements of the heart. The vast majority of men are only sensitive to five or six passions, in the circle of which their lives are passed and which define the boundaries of their imaginations. Take away love and hate, pleasure and pain, hope and fear, and they will feel nothing.

But persons of a nobler character can be moved in thousands of different ways. It seems that they can receive ideas and sensations which surpass the ordinary norms of nature.
\flright{\textsc{Abbe Prevost}, \emph{Manon Lescaut}}
\end{quotex}


Prevost was not far off. Psychologists have identified six or so universal emotions\footnote{\url{https://people.ece.cornell.edu/land/OldStudentProjects/cs490-95to96/HJKIM/emotions.html}} with their corresponding facial expressions: disgust, sadness, happiness, fear, anger, surprise. Therefore, the development of the emotional centre in order to experience more subtle states is an important part of esoteric training.

The intellectual centre, on the other hand, is completely undeveloped and therefore must be trained. The result is that the child learns programmed responses so his deeper self gets submerged. Boris Mouravieff compares this programming to recordings. Dialog becomes rote and automatic. He writes:

\begin{quotex}
This procedure is mechanical. We can easily observe this in any conversation between a number of persons who know each other slightly. Such an interchange necessarily falls to an elementary level of the most banal interests: weather, political news, or the city. We hear these records being played, turning continually and passing from one person to another, each with their faces congealed in a grimace which---we commonly agree --- gives evidence of an amiable attitude.
\end{quotex}


He continues:

\begin{quotex}
The recording continues practically forever, as the disc library is vast and the recording apparatus very sensitive. When a person speaks, it is generally easy to distinguish whether his records are playing or whether he speaks from some deeper part of himself. ... One moment it is `I' who speaks then, unnoticed, it is no longer I; a recording from the past begins to play in me. A curious thing: once a record has been started, it is almost impossible to stop it before it has run through its content.
\end{quotex}


Therefore popular fiction needs to be limited to Prevost's passions and the universal emotions, otherwise it will be misunderstood. Those of a nobler character are simply bored by popular culture due to its self-imposed limitations.

Creators of popular culture are certainly aware of such limitations, so they create for that level. Of course, in most cases they may simply be incapable of creating art and literature of High Culture. If you have ever taken acting or screenwriting classes, as I have, you will start looking at shows and movies in a much different way. Instead of being captivated by the spectacle, you will perceive ``behind the scenes'', to see how the director or writer arrange things to create a purposeful effect. That effect is intended to appeal to certain inner states of the audience or invoke predetermined reactions. That makes the viewer a purely passive receptacle of the ideas of others. In High Culture, effort is required to participate in the creation of the author.

\paragraph{Psychological Manipulation}


Generally speaking, the writers are more intelligent and better educated than the popular audience. If you listen to common everyday speech about personal events, you will notice the meanderings, lack of focus, and inability to get to the point. On the other hand, experienced writers follow a plan. For example, Aristotle wrote about the structure and plots of plays. Others followed, such as the Romans Horace and Aelius Donatus.

Writers will also employ insights from psychology, particularly the excitation of the eros and thumos drives. That is popularly called ``sex and violence''. Only High Culture will try to engage the nous. So creators of popular culture will stick to the themes mentioned by Abbe Provost and the more primitive, universally recognized, emotions as described above.

More recently, writers are learning to appeal directly to the subconscious. Specifically, the structure of myths described by \textbf{Joseph Campbell} in \emph{Hero With A Thousand Faces} is employed by screenwriters. More recently, \emph{The Origins and History of Consciousness} by Jung's disciple \textbf{Erich Neumann} is serving that purpose. Of course, both works have similarities due to their common source in \textbf{Carl Jung}. In this way, the writer is able to reach into the subconscious archetypes of the audience. For more information, check out \emph{Consciousness, Personality \& Hero Sequence The Origin and History of Consciousness: A New Structure for Stories?}\footnote{\url{https://greathousestories.files.wordpress.com/2015/12/origin-history-of-consciousness1.pdf}} by a professional writer.

\paragraph{Jack Ryan}


As an example, I will use the Jack Ryan show that was streamed on Amazon. I did not know at the time that this character has appeared in several books and movies. However, in postmodern fashion, in which ``there is nothing but the text'', I'll treat the series as self-contained, just as I watched it. I'll try to relate the show to Joseph Campbell's hero sequence as described in the link above. The sequence is incomplete and some stages may be out of order, probably because the story has to appeal to the imagination of typical viewers.

\textbf{The Ordinary World}. Jack Ryan seems to be an ordinary cubicle rat, albeit with some extraordinary credentials. He has an advanced degree from Boston College and was a financial analyst before his CIA job. Of course, for popular TV it is essential to appeal to the Walter Mitty fantasies\footnote{\url{https://en.wikipedia.org/wiki/Walter\_Mitty}} of the audience. Thus, Ryan is in demand by rich and powerful men, yet he prefers life in his cubicle. He also has six pack abs, which is unlikely from sitting in front of a computer all day.

\textbf{Love Interest}. This is not in the screenwriter's or Campbell's list, although it is needed in most any story, from Homer to Virgil and beyond. In medieval romances, the Lady would set the task for the Knight. Not in this case, since Cathy is opposed to Ryan's secret life. Cathy is an attractive physician, not a waitress from Denny's -- nothing but the best for Ryan. She turns out to be a cheap date. He doesn't even have to worry about the \#Metoo movement, since Cathy initiates their first sexual encounter. Even better for him, she insists on keeping their relationship casual, as though that would disappoint him. Au contraire, every guy would love to hear that, since he interprets it as a ``friend with benefits''.

\textbf{Call to Adventure}. He notices some abnormal transactions which he believes can be traced to a terrorist named Suleiman. He tries to bring his theory up the command chain but experiences resistance.

\textbf{Meeting with the Mentor}. Greer is Jack's cynical and disillusioned mentor who initially doesn't accept Ryan's conclusion, but then guides him through the process. He is also the ``good Muslim'', unlike Suleiman.

\textbf{Crossing the Threshold}. The honchos eventually figure it out. In spectacular fashion the send a black helicopter to pick up Ryan while he is at an elegant party. That impresses his new girlfriend Cathy, again, another plebeian fantasy. Ryan is whisked away to the Near East with ``unfamiliar rules and values''.

\textbf{Test, Allies and Enemies}. The hero Ryan is tested in unique ways and is forced to distinguish between allies and enemies. Even his big bosses have doubts and prefer to ``play it safe''. That is a test so that Ryan is forced to stand his ground. Suleiman's wife betrays her husband and turns into an ally. Her motivation is to ``protect the children'', as if you had never heard that before. That creates another test: Ryan has agreed to keep the children out of harm's way.

\textbf{The Approach}. The hero Ryan and his allies have to prepare for an unknown attack by Suleiman.

\textbf{The Ordeal and the Road Back}. Ultimately Ryan is alone with Suleiman in a running gun battle through the subway system. Suddenly the desk jockey is a crack shot, even though breathless and anxious. In the face of death, he becomes more than a nerd. He is now the hero, and saved the King (or President) no less.

\textbf{The Resurrection}. The final test is to bring Suleiman's wife and children to safety. Ryan is a man of his word and the conflict is over.

\textbf{Return with the Elixer}. Ryan returns home as a hero and is offered a promotion. By his facial expressions, Ryan is not happy with another desk job. Meanwhile, his mentor is transferred to the Moscow office. By leaving him a ticket, the presumption is that Ryan will join him there for another adventure. He's better off leaving his whiny blonde girlfriend. In Moscow, the Slavic women are more exciting, even if somewhat treacherous.

\paragraph{Conclusion}


A government study has shown\footnote{\url{https://www.bloomberg.com/amp/news/articles/2018-12-10/screen-time-changes-structure-of-kids-brains-60-minutes-says}} that screen time has deleterious effects on brain development in young children. A similar study needs to be done on adults, many of whom watch propaganda all day. This revelation should convince you about deliberate psychological manipulation in the media. Maybe it is done just to make a buck, but in many shows, it goes well beyond that.


\flrightit{Posted on 2018-12-11 by Cologero}

\begin{center}* * *\end{center}

\begin{footnotesize}
\begin{sffamily}
\texttt{swagmoto on 2018-12-12 at 20:45 said:}

\url{https://www.youtube.com/watch?v=tM8nR9sJbMQ}


\hfill

\texttt{Tom on 2018-12-13 at 12:11 said:}

A very revealing example of this, in which the ``magician reveals his tricks'', is Christopher Nolan's ``Inception''. The protagonists of the film are experts at creating and manipulating dream worlds. Their mission is to put to sleep \& surreptitiously invade the dream of an heir to a large company and plant the idea in him that he should disinherit himself (make your own connections to the political history of the late 20th century). They have to go deep into his subconscious within the dream to be able to do this, and then fabricate a narrative for him that takes account of his personal factors and strained relationship with his father in order to bring all of these internal tensions to a particular resolution that ends in him adopting they idea they want to implant in him.

This is a pattern of ``black magic'' (the goal being to make the target less free; subjected to the will of the magician) and is exactly how modern television and film today are set up -- to put the target into a state of reduced consciousness with images and appeals to eros, etc., and then to start working through the stages of human psychology in order to ultimately implant the idea of the filmmaker, all while leaving the viewer totally unaware that he has been psychologically ``raped''.

What's new about this isn't so much the reliance on traditional mythological structures (the ancients obviously used similar knowledge to create stories), but the power of modern cinema to suppress the reasoning faculty through imagery \& sound and more effectively put the audience into a ``dream'' state. This way, ideas can be implanted that the audience would otherwise reject out of hand. The lack of a religious foundation or any true development of intellectuality makes modern man particularly defenseless against these kinds of suggestion.

There are nevertheless grey areas --- opera, for example. Opera employs spectacle and sound in order to bring the audience into a desired psychological state, and uses traditional mythological structures in order to communicate ideas, but most would agree that going to see a traditionally staged \& performed opera is a far more uplifting experience than watching the latest superhero film. It requires some intellectual and emotional participation from the audience, but generally not as much as, say, Shakespeare. Even a relatively uncultured person can be touched by Richard Wagner.

This subject of the proper employment of art, music, and storytelling is discussed at length in Book III of the Republic, which I will quote here:

``And therefore, I said, Glaucon, musical training is a more potent instrument than any other, because rhythm and harmony find their way into the inward places of the soul, on which they mightily fasten, imparting grace, and making the soul of him who is rightly educated graceful, or of him who is ill-educated ungraceful; and also because he who has received this true education of the inner being will most shrewdly perceive omissions or faults in art and nature, and with a true taste, while he praises and rejoices over and receives into his soul the good, and becomes noble and good, he will justly blame and hate the bad, now in the days of his youth, even before he is able to know the reason why; and when reason comes he will recognise and salute the friend with whom his education has made him long familiar. ``

\hfill

\texttt{Cologero on 2018-12-14 at 16:23 said:}

@Tom, Good points

As you point out, the goal of ``black magic'' is to make the subject less free, rather than to increase his freedom.

Given what you wrote, I think we can make a distinction between a mythos that arises organically and one that is deliberately constructed.

For example, a good opera might use a myth that arose from the spiritual formation of a folk. In that case, I would think the effect is salutary.

On the other hand, if the structure of myth is deliberately employed for commercial ends, the effect is most likely destructive.

\end{sffamily}
\end{footnotesize}

